\chapter{System Design}

\section{System Architecture}

The Indian Sign Language Detection System follows a modular architecture with the following main components:

\begin{itemize}
    \item \textbf{Input Module}: Video capture and preprocessing
    \item \textbf{Detection Module}: Hand landmark detection
    \item \textbf{Classification Module}: Gesture recognition
    \item \textbf{Output Module}: Text and speech generation
    \item \textbf{User Interface}: Web dashboard
    \item \textbf{Data Management}: Training data and model storage
\end{itemize}

\section{Basic Modules}

\begin{table}[H]
\centering
\begin{tabular}{|c|c|c|}
\hline
\textbf{Module} & \textbf{Input} & \textbf{Output} \\
\hline
Hand Detection & RGB frame stream & 21-point hand landmarks \\
\hline
Gesture Classification & Normalized landmarks & Predicted gesture label \\
\hline
Point History Tracking & Fingertip trajectory sequence & Motion-gesture class \\
\hline
Text-to-Speech & Gesture text/token & Audible speech output \\
\hline
Dashboard Interface & User controls and camera feed & Training/testing interaction \\
\hline
\end{tabular}
\caption{Basic Modules and I/O Mapping}
\end{table}

\subsection{Hand Detection Module (MediaPipe)}

\begin{itemize}
    \item Input: RGB video frame (640x480 resolution)
    \item Process: Neural network-based hand detection
    \item Output: 21 hand landmarks with x, y, z coordinates
    \item Handles: Single and multiple hand detection, occlusion robustness
\end{itemize}

\subsection{Gesture Classification Module}

\begin{itemize}
    \item Input: Preprocessed landmark coordinates (84 values = 21 landmarks × 4)
    \item Process: CNN/Dense network classification
    \item Output: Gesture class (recognizable ISL gestures)
    \item Trained on: Custom dataset (100+ images per gesture)
\end{itemize}

\subsection{Point History Tracking Module}

\begin{itemize}
    \item Input: Sequence of finger tip positions
    \item Process: Historical trajectory analysis
    \item Output: Complex gesture classification (swipe, circle, etc.)
    \item Buffer: Stores last 16 frames of trajectory data
\end{itemize}

\subsection{Text-to-Speech Module}

\begin{itemize}
    \item Input: Recognized gesture text or user-typed text
    \item Process: pyttsx3 engine synthesis
    \item Output: Audio file playback
    \item Threading: Non-blocking speech generation
\end{itemize}

\section{Data Flow Diagram}

\begin{figure}[H]
\centering
\fbox{
\begin{minipage}{0.9\textwidth}
\includegraphics[width=\textwidth]{diagrams/dfd_placeholder.png}
\end{minipage}
}
\caption{Data Flow Diagram - System processes from input to output}
\label{fig:dfd}
\end{figure}

\subsection{Detailed Diagram Description (For Generation)}

Use this structure while generating the DFD:
\begin{itemize}
    \item \textbf{External Entity}: User (shows gesture to camera and receives text/speech output)
    \item \textbf{Input Process}: Webcam Capture (reads real-time frames)
    \item \textbf{Process 1}: Frame Preprocessing (flip, resize, color conversion BGR to RGB)
    \item \textbf{Process 2}: Hand Landmark Detection (MediaPipe extracts 21 points per hand)
    \item \textbf{Process 3}: Landmark Preprocessing (normalization, feature vector creation)
    \item \textbf{Process 4}: Keypoint Classification (predict static gesture label)
    \item \textbf{Process 5}: Point History Classification (predict motion gesture from trajectory)
    \item \textbf{Data Store D1}: Training Dataset (gesture images and CSV features)
    \item \textbf{Data Store D2}: Model Files (.hdf5/.tflite classifiers and label CSVs)
    \item \textbf{Process 6}: Output Formatter (map class ID to gesture text)
    \item \textbf{Process 7}: Text-to-Speech Engine (convert text to audio)
    \item \textbf{Output}: Dashboard/Window display + speaker audio
    \item \textbf{Main Flow}: User -> Webcam -> Preprocessing -> Detection -> Preprocessing -> Classification -> Text -> Speech/Display -> User
\end{itemize}

The data flow follows this sequence:
\begin{enumerate}
    \item User provides video input (webcam feed)
    \item Hand landmarks are detected using MediaPipe
    \item Landmarks are preprocessed and normalized
    \item Gesture classifier processes the landmarks
    \item Point history classifier analyzes trajectory
    \item Gesture name is generated
    \item Text-to-speech converts output to audio
    \item Results displayed in UI (text + speech)
\end{enumerate}

\section{Entity-Relationship Diagram}

\begin{figure}[H]
\centering
\fbox{
\begin{minipage}{0.9\textwidth}
\includegraphics[width=\textwidth]{diagrams/er_diagram.png}
\end{minipage}
}
\caption{ER Diagram - Database and data structure relationships}
\label{fig:er}
\end{figure}

\subsection{Detailed Diagram Description (For Generation)}

Use the following entities, keys, and relationships:
\begin{itemize}
    \item \textbf{Entity: Gesture}
    \begin{itemize}
        \item Attributes: gesture\_id (PK), gesture\_name, created\_date
    \end{itemize}
    \item \textbf{Entity: TrainingImage}
    \begin{itemize}
        \item Attributes: image\_id (PK), gesture\_id (FK), image\_path, timestamp
    \end{itemize}
    \item \textbf{Entity: Model}
    \begin{itemize}
        \item Attributes: model\_id (PK), model\_type, accuracy, trained\_date
    \end{itemize}
    \item \textbf{Entity: Session}
    \begin{itemize}
        \item Attributes: session\_id (PK), start\_time, end\_time, recognized\_gestures
    \end{itemize}
    \item \textbf{Entity: PredictionLog}
    \begin{itemize}
        \item Attributes: log\_id (PK), session\_id (FK), gesture\_id (FK), confidence, log\_time
    \end{itemize}
\end{itemize}

Relationship guidance:
\begin{itemize}
    \item Gesture (1) to TrainingImage (N)
    \item Session (1) to PredictionLog (N)
    \item Gesture (1) to PredictionLog (N)
    \item Model can be shown as independent configuration entity used by classification pipeline
\end{itemize}

\subsection{Main Entities}

\begin{table}[H]
\centering
\begin{tabular}{|c|c|c|}
\hline
\textbf{Entity} & \textbf{Attributes} & \textbf{Description} \\
\hline
\multirow{3}{*}{Gesture} & gesture\_id (PK) & Unique gesture identifier \\
& gesture\_name & Name of the gesture \\
& created\_date & Creation timestamp \\
\hline
\multirow{4}{*}{TrainingImage} & image\_id (PK) & Unique image identifier \\
& gesture\_id (FK) & Reference to gesture \\
& image\_path & Location of image file \\
& timestamp & Collection time \\
\hline
\multirow{3}{*}{Model} & model\_id (PK) & Unique model identifier \\
& model\_type & (keypoint/history) \\
& accuracy & Validation accuracy \\
\hline
\multirow{3}{*}{Session} & session\_id (PK) & Unique session identifier \\
& recognized\_gestures & List of gestures \\
& timestamp & Session start time \\
\hline
\end{tabular}
\caption{Main Database Entities}
\end{table}

\section{Use Case Diagram}

\begin{figure}[H]
\centering
\fbox{
\begin{minipage}{0.9\textwidth}
\includegraphics[width=\textwidth]{diagrams/usecase_diagram.png}
\end{minipage}
}
\caption{Use Case Diagram - System interactions with actors}
\label{fig:usecase}
\end{figure}

\subsection{Detailed Diagram Description (For Generation)}

Use Case diagram specification:
\begin{itemize}
    \item \textbf{Primary Actor}: User
    \item \textbf{Secondary Actor}: Administrator/Developer (optional for model management)
    \item \textbf{System Boundary}: ISL Detection System
    \item \textbf{Use Cases}:
    \begin{itemize}
        \item Train New Gesture
        \item Collect Training Images
        \item Train/Update Model
        \item Test Gesture Recognition
        \item Recognize Gesture in Real Time
        \item View Training Data Statistics
        \item Type Using Gestures
        \item Convert Text to Speech
        \item Manage Gesture Mapping
    \end{itemize}
\end{itemize}

Relation guidance:
\begin{itemize}
    \item \texttt{Test Gesture Recognition} includes \texttt{Recognize Gesture in Real Time}
    \item \texttt{Type Using Gestures} includes \texttt{Recognize Gesture in Real Time}
    \item \texttt{Type Using Gestures} includes \texttt{Convert Text to Speech}
    \item \texttt{Train New Gesture} includes \texttt{Collect Training Images} and \texttt{Train/Update Model}
\end{itemize}

Main use cases:
\begin{enumerate}
    \item Train New Gesture
    \item Test Gesture Recognition
    \item Type Using Gestures
    \item View Training Data
    \item Generate Speech Output
    \item Manage Models
\end{enumerate}

\section{Class Diagram}

\begin{figure}[H]
\centering
\fbox{
\begin{minipage}{0.9\textwidth}
\includegraphics[width=\textwidth]{diagrams/class_diagram.png}
\end{minipage}
}
\caption{Class Diagram - Object-oriented system structure}
\label{fig:class}
\end{figure}

\subsection{Detailed Diagram Description (For Generation)}

Class diagram structure to generate:
\begin{itemize}
    \item \textbf{SimpleHandDetector}
    \begin{itemize}
        \item Methods: detect\_hands(frame), get\_landmarks(), get\_handedness()
    \end{itemize}
    \item \textbf{KeyPointClassifier}
    \begin{itemize}
        \item Methods: \_\_init\_\_(model\_path), \_\_call\_\_(landmark\_list)
    \end{itemize}
    \item \textbf{PointHistoryClassifier}
    \begin{itemize}
        \item Methods: \_\_init\_\_(model\_path), \_\_call\_\_(history\_vector)
    \end{itemize}
    \item \textbf{SimpleGestureRecognizer}
    \begin{itemize}
        \item Methods: preprocess\_landmarks(), predict\_gesture(), map\_label()
    \end{itemize}
    \item \textbf{SpeechEngine}
    \begin{itemize}
        \item Methods: speak\_text(text), \_speak\_worker()
    \end{itemize}
    \item \textbf{DashboardController}
    \begin{itemize}
        \item Methods: train\_mode(), test\_mode(), typing\_mode(), view\_data\_mode()
    \end{itemize}
\end{itemize}

Association guidance:
\begin{itemize}
    \item DashboardController uses SimpleHandDetector, KeyPointClassifier, PointHistoryClassifier, and SpeechEngine
    \item SimpleGestureRecognizer depends on KeyPointClassifier and label mappings
\end{itemize}

\subsection{Key Classes}

\begin{itemize}
    \item \textbf{KeyPointClassifier}: Gesture classification model
    \item \textbf{PointHistoryClassifier}: Complex gesture recognition
    \item \textbf{SimpleGestureRecognizer}: High-level gesture matching
    \item \textbf{SimpleHandDetector}: Hand detection wrapper
    \item \textbf{SpeechEngine}: Text-to-speech functionality
\end{itemize}

\section{Sequence Diagram}

\begin{figure}[H]
\centering
\fbox{
\begin{minipage}{0.9\textwidth}
\includegraphics[width=\textwidth]{diagrams/sequence_diagram.png}
\end{minipage}
}
\caption{Sequence Diagram - Interaction timeline for gesture recognition}
\label{fig:sequence}
\end{figure}

\subsection{Detailed Diagram Description (For Generation)}

Use the following participants and message order:
\begin{itemize}
    \item Participants: User, Camera, Preprocessor, HandDetector, LandmarkPreprocessor, KeypointClassifier, HistoryClassifier, OutputManager, SpeechEngine, UI
    \item Sequence:
    \begin{enumerate}
        \item User -> Camera: show gesture
        \item Camera -> Preprocessor: raw frame
        \item Preprocessor -> HandDetector: processed frame
        \item HandDetector -> LandmarkPreprocessor: hand landmarks
        \item LandmarkPreprocessor -> KeypointClassifier: normalized vector
        \item LandmarkPreprocessor -> HistoryClassifier: trajectory vector (optional)
        \item KeypointClassifier -> OutputManager: class id + confidence
        \item HistoryClassifier -> OutputManager: motion class (optional)
        \item OutputManager -> UI: display text label and confidence
        \item OutputManager -> SpeechEngine: speak(label)
        \item SpeechEngine -> User: audio output
    \end{enumerate}
\end{itemize}

\subsection{Gesture Recognition Sequence}

\begin{enumerate}
    \item User shows gesture to camera
    \item System captures video frame
    \item MediaPipe detects hand landmarks
    \item Landmarks are normalized and preprocessed
    \item Classifier predicts gesture type
    \item If complex gesture: analyze point history
    \item Generate text output
    \item Trigger text-to-speech
    \item Display results in UI
\end{enumerate}

\section{Activity Diagram}

\begin{figure}[H]
\centering
\fbox{
\begin{minipage}{0.9\textwidth}
\includegraphics[width=\textwidth]{diagrams/activity_diagram.png}
\end{minipage}
}
\caption{Activity Diagram - Processing workflow and decision points}
\label{fig:activity}
\end{figure}

\subsection{Detailed Diagram Description (For Generation)}

Activity flow to generate:
\begin{enumerate}
    \item Start
    \item Initialize camera and models
    \item Capture frame
    \item Preprocess frame
    \item Detect hand landmarks
    \item Decision: Hand detected?
    \begin{itemize}
        \item No -> Capture next frame
        \item Yes -> Continue
    \end{itemize}
    \item Preprocess landmarks
    \item Predict static gesture
    \item Update point history
    \item Decision: Enough history for motion classification?
    \begin{itemize}
        \item No -> Skip motion class
        \item Yes -> Predict motion gesture
    \end{itemize}
    \item Map class to text label
    \item Display result in UI
    \item Decision: Auto-speak enabled?
    \begin{itemize}
        \item Yes -> Send text to speech engine
        \item No -> Continue
    \end{itemize}
    \item Decision: Exit key/button pressed?
    \begin{itemize}
        \item No -> Loop to capture frame
        \item Yes -> Stop camera, release resources, End
    \end{itemize}
\end{enumerate}

\subsection{Main Activities}

\begin{enumerate}
    \item Video Capture: Initialize camera and capture frames
    \item Hand Detection: Apply MediaPipe landmark detection
    \item Preprocessing: Normalize and scale landmark data
    \item Classification: Feed to gesture classifier
    \item Decision: Confidence threshold check
    \item Output Generation: Create text and speech
    \item Display: Show results to user
\end{enumerate}

\section{Data Design Summary}

\begin{table}[H]
\centering
\begin{tabular}{|c|c|c|}
\hline
\textbf{Data Component} & \textbf{Storage Form} & \textbf{Design Purpose} \\
\hline
Training images & Gesture-wise folder structure & Label organization and scalable data collection \\
\hline
Landmark features & CSV rows (numeric vectors) & Lightweight model-training input format \\
\hline
Label mapping & CSV label files & Stable class-index to gesture-name mapping \\
\hline
Configuration & JSON metadata file & Persisted gesture and dashboard settings \\
\hline
Trained models & HDF5 and TFLite files & Training-time and deployment-time compatibility \\
\hline
\end{tabular}
\caption{Data Design Overview}
\end{table}

\section{Procedural Design Summary}

\begin{table}[H]
\centering
\begin{tabular}{|c|p{0.62\textwidth}|}
\hline
\textbf{Procedure} & \textbf{Algorithmic Logic} \\
\hline
Landmark preprocessing & Wrist-centered normalization and flattening to create translation/scale invariant features. \\
\hline
Gesture inference & Pass normalized vector to classifier and select highest confidence class. \\
\hline
Temporal motion handling & Maintain fixed-length history window and classify trajectory patterns. \\
\hline
Speech generation & Run non-blocking speech synthesis thread for recognized tokens/text. \\
\hline
\end{tabular}
\caption{Procedural Design and Algorithm Mapping}
\end{table}

\section{Design Patterns and Principles}

\subsection{Architectural Patterns}

\subsubsection{Pipeline Pattern}

The system implements a linear pipeline pattern:
\begin{itemize}
    \item Input (webcam frame) → Preprocessing → Detection → Classification → Output
    \item Each stage is independent, allowing for modular testing and replacement
    \item Data flows sequentially without feedback loops
    \item Exceptions are handled at each stage with graceful fallback
\end{itemize}

\subsubsection{MVC (Model-View-Controller) in Dashboard}

The Streamlit dashboard implements MVC principles:
\begin{itemize}
    \item \textbf{Model}: Trained gesture classifiers and data structures
    \item \textbf{View}: Streamlit UI components and visualizations
    \item \textbf{Controller}: Mode functions (train\_mode, test\_mode, typing\_mode)
    \item Separation ensures UI logic doesn't interfere with core recognition
\end{itemize}

\subsection{Design Principles}

\subsubsection{Separation of Concerns}

\begin{itemize}
    \item \textbf{Detection Logic}: Isolated in MediaPipe wrapper class
    \item \textbf{Classification Logic}: Separate modules for keypoint and history classifiers
    \item \textbf{UI Logic}: Contained in Streamlit dashboard
    \item \textbf{Output Logic}: Separate handler for text and speech output
    \item Each component has single responsibility with clear interfaces
\end{itemize}

\subsubsection{SOLID Principles}

\begin{itemize}
    \item \textbf{Single Responsibility}: Each class has one reason to change
    \item \textbf{Open/Closed}: System open for extension (new gestures) but closed for modification
    \item \textbf{Liskov Substitution}: Classifiers have consistent interface for substitution
    \item \textbf{Interface Segregation}: Classes expose only necessary methods to other components
\end{itemize}

\section{Error Handling and Exception Strategy}

\subsection{Exception Handling}

\begin{itemize}
    \item \textbf{Camera Integration}: Graceful handling of camera unavailability with retry logic
    \item \textbf{Model Loading}: Validation of model files with fallback to default models
    \item \textbf{Hand Detection Failures}: Continuous monitoring with "No hand" feedback
    \item \textbf{Classification Errors}: Confidence thresholding to identify uncertain predictions
    \item \textbf{Audio System}: Fallback to visual-only output if text-to-speech unavailable
\end{itemize}

\section{Performance and Scalability}

\subsection{Performance Optimizations}

\begin{itemize}
    \item \textbf{Model Quantization}: TFLite reduces model by 60\% with minimal accuracy loss
    \item \textbf{Vectorized Operations}: NumPy for fast preprocessing instead of Python loops
    \item \textbf{Threading}: Background thread for speech synthesis prevents UI blocking
    \item \textbf{Memory Efficiency}: Reusable buffers for frame storage
\end{itemize}

\subsection{Scalability}

\begin{itemize}
    \item \textbf{Gesture Vocabulary}: Scales to 20-30 practical gestures; training time ~30 min per gesture
    \item \textbf{Model Size}: Remains <1 MB even with 30 gestures
    \item \textbf{Deployment}: Containerization (Docker) feasible for cloud deployment
    \item \textbf{Multi-Instance}: Load balancing possible for web service variants
\end{itemize}

\section{Configuration and Customization}

\subsection{Tunable Parameters}

\begin{itemize}
    \item \textbf{Camera Resolution}: Default 640x480, adjustable for different hardware
    \item \textbf{Detection Confidence}: 0.5-0.7 tunable threshold
    \item \textbf{Classification Confidence}: Minimum 0.7 for reporting gestures
    \item \textbf{Speech Rate}: 150-200 WPM configurable
    \item \textbf{History Buffer}: 16 frames for temporal analysis
\end{itemize}

\subsection{Extensibility Points}

\begin{itemize}
    \item New gesture models: Add TFLite file + label CSV
    \item Alternative detection: Replace MediaPipe with alternative library
    \item Custom classifiers: Plug in alternative ML models
    \item Language support: Add language-specific labels and TTS models
\end{itemize}

\section{Security Design}

\subsection{Data Protection}

\begin{itemize}
    \item \textbf{Local Processing}: All data stays on user's machine; no cloud upload
    \item \textbf{No Persistent Logging}: Video frames and audio not stored permanently
    \item \textbf{User Control}: Training data in user-controlled directories
    \item \textbf{Input Validation}: Frame dimensions, landmark confidences verified before processing
\end{itemize}

\section{Database and Data Persistence}

\subsection{Storage Strategy}

\begin{itemize}
    \item \textbf{Training Data}: Organized in folder hierarchy by gesture name
    \item \textbf{Models}: Stored as HDF5 (training) and TFLite (deployment) files
    \item \textbf{Labels}: CSV format for class-index to name mapping
    \item \textbf{Logs}: CSV-based session and prediction logs
    \item \textbf{Advantages}: Simple, human-readable, version-control friendly
\end{itemize}

\section{User Interface and Interaction}

\subsection{Interface Layers}

\begin{itemize}
    \item \textbf{Web Dashboard} (primary): Streamlit with training/testing/typing modes
    \item \textbf{Real-Time Feedback}: Live video, landmarks, and confidence display
    \item \textbf{Settings}: User-adjustable parameters (speech rate, confidence threshold)
    \item \textbf{Visualization}: Gesture history, training statistics, model performance charts
\end{itemize}

\subsection{Key Interaction Flows}

\subsubsection{Training Workflow}
\begin{enumerate}
    \item User selects "Train New Gesture" mode
    \item Specifies gesture name and image count target
    \item Captures images via webcam in dashboard
    \item Images automatically labeled and stored
    \item Model retraining automatically triggered
    \item System ready with new gesture integrated
\end{enumerate}

\subsubsection{Recognition Workflow}
\begin{enumerate}
    \item User selects "Test & Recognize" mode
    \item Real-time video shows hand detection
    \item Detected hand landmarks overlaid on video
    \item Recognized gesture displayed with confidence
    \item Gesture automatically spoken aloud
    \item History maintained for reference
\end{enumerate}

\section{Design Validation Points}

The system design has been validated against the following criteria:

\begin{itemize}
    \item \textbf{Modularity}: Each component can be tested, replaced, or upgraded independently
    \item \textbf{Maintainability}: Code follows clear structure with minimal technical debt
    \item \textbf{Extensibility}: New gestures, models, and features can be added without major refactoring
    \item \textbf{Robustness}: Error handling covers primary failure modes
    \item \textbf{Performance}: Meets real-time processing requirements (\textgreater 20 FPS)
    \item \textbf{Usability}: UI designed for non-technical end-users
    \item \textbf{Scalability}: Architecture supports gesture vocabulary growth and multi-user deployment
\end{itemize}
